% ============================================================
% Chapter 8 — Measurement Methodology
% ============================================================
\chapter{Measurement Methodology}
\label{ch:methodology}

\roadmap{This chapter specifies how the platform's data support the
project's measurement claims. It presents the cross-modal convergent
validation design relating facial-expression engagement estimates to
gaze-on-AOI dwell; the nine-construct executive-function text-mining
pipeline and its calibration metadata; the self-report probe that anchors
validation; the pre/post learning-gain method; a catalogue of analyses
the collected data can support; and a candid statistics-readiness
assessment for the proposal's required tests. Throughout it names the
limits of each design rather than deferring them to
\cref{ch:limitations}.}

\section{Cross-modal convergent validation}
\label{sec:methodology-validation}

RQ1 asks whether browser-based sensing can detect learner states with
useful fidelity. Absent live human coding, the design validates one
machine channel against another: facial-expression-derived engagement
against gaze-on-AOI dwell as a \emph{behavioral reference criterion}. This
is convergent validation in the sense of \citet{campbell1959}---agreement
between two methods targeting the same construct---and explicitly
\emph{not} criterion validity against ground truth. Both channels can err,
and their errors may correlate; the design establishes convergence, and
the report never overstates it.

\paragraph{Windowing.} Analyses use fixed windows aligned to the
sync-anchor time base---either the server's 30\,s window (10\,s stride)
for comparison with \dbtable{AffectiveStateWindow}, or a finer 10\,s
window for resolution---joined to the CSV time grid
(\cref{sec:sensing-export}).

\paragraph{Facial features.} The engagement estimate is taken either as
the server-derived engagement score or as a function of the eighteen
action units aggregated per window (mean intensity), with AU06/AU12
(positive affect), AU04 (concentration or negative affect), and AU01/AU02
(surprise/attention) as the interpretable candidates
\citep{grafsgaard2013}.

\paragraph{On-task definition.} Per epoch type, on-task is defined as gaze
dwelling in the epoch's expected AOI bucket: the lesson/PDF bucket during
reading, the chatbot bucket during dialogue, and the active panel during
an intervention. The continuous per-epoch alignment of
\cref{eq:alignment} provides a graded on-task index; a median split yields
a binary on-task label.

\paragraph{Agreement statistics.} With engagement and on-task computed per
window, the design supports (i) a point-biserial or Pearson correlation
between continuous engagement and continuous alignment, (ii) Cohen's
$\kappa$ \citep{cohen1960} between the two binarized signals across
windows, and (iii) a confusion matrix with sensitivity and specificity.
These are reported per participant and pooled in \cref{ch:results}.

\paragraph{Threats.} Webcam gaze is region-level and noisy; the AOI
expected weights are placeholders; and both channels share the webcam, so
face-loss and poor lighting degrade them together. The design therefore
measures convergence under shared conditions, and the discussion
(\cref{ch:discussion}) treats disagreement as informative about channel
limits, not merely as error.

\section{Executive-function text mining}
\label{sec:methodology-ef}

In parallel to the facial channel, an LLM text-mining pipeline detects
psychologically motivated constructs in learner utterances, giving affect
and cognition a linguistic detection channel
(\cref{sec:bg-ef}).\footnote{\code{apps/api/src/text-mining}} Nine
constructs are configured, seeded from code into versioned prompt records
and overridable per course; \cref{tab:ef-constructs} lists them with their
label type and the implementation's own feasibility rating.

\begin{table}[htbp]
\centering
\small
\caption{The nine text-mining constructs, with label type and the
implementation's feasibility rating (5 = most feasible).}
\label{tab:ef-constructs}
\begin{tabular}{@{}p{0.34\linewidth}p{0.28\linewidth}c@{}}
\toprule
Construct & Label type & Feasibility \\
\midrule
Metacognition (general) & binary & 5 \\
Metacognitive monitoring & binary & 5 \\
Attention / mind-wandering & binary & 2 \\
Working memory load & ordinal (low/med/high) & 2 \\
Cognitive flexibility & binary & 1 \\
Confusion & binary (+severity 1--5) & 5 \\
Frustration & binary & 4 \\
Engagement / off-task & on-task/off-task & 4 \\
Boredom & binary & 2 \\
\bottomrule
\end{tabular}
\end{table}

\paragraph{Flow.} Detection fires per learner utterance, fire-and-forget,
from both dialogue and chatbot turns. Each enabled construct is one LLM
call under a bounded-concurrency scheduler; results are written to
\dbtable{EfDetection}. The provider and model default to the learner's
configured LLM (typically a small GPT-class model) and can be overridden
by teacher settings.

\paragraph{Calibration metadata.} Every detection stores the fields needed
for calibration and audit: a confidence, an optional severity (populated
for confusion), a free-text rationale, an optional warning (for the two
low-feasibility EF constructs, which the prompts themselves flag as
unvalidated), the prompt version, provider, model, and latency. This lets
later analysis condition on confidence, exclude low-feasibility constructs,
and track prompt drift. One caveat: failed detections are retained as
error rows and must be filtered before analysis.

\section{Self-report instrument}
\label{sec:methodology-selfreport}

The study's affect criterion is a mandatory self-report
probe.\footnote{\code{apps/web/src/components/EmotionSurveyModal.tsx};
\code{ActivityAction.EMOTION_SELF_REPORT}} Every fifteen minutes a student
is shown a non-dismissable modal offering five options---engaged, bored,
confused, frustrated, neutral---and the choice is written to
\dbtable{ActivityLog} metadata. The five options were chosen to match the
machine-detected state space almost exactly, which is what makes the
validation of item 15 possible. Two caveats are stated up front: the probe
is sparse (one report per fifteen minutes) relative to the dense state
windows, so each report maps to many windows; and it was added on
2026-07-01, so approximately 54 earlier sessions carry no self-report and
are unusable for item-15 validation. \todo{Confirm the count of
post-2026-07-01 sessions with self-report data.}

\section{Learning gain}
\label{sec:methodology-gain}

RQ2's learning outcome is measured by a pre/post design. Two parallel-form
assessments were seeded on the same course---a pre-study assessment and a
post-study assessment with identical topic coverage, structure, point
weights, and difficulty tags but entirely new
items.\footnote{\code{apps/api/prisma/scripts/seed-cs601-assessment.ts};
\code{seed-cs601-post-assessment.ts}} A learner's pre and post scores are
obtained by summing \dbtable{GradingResult.score} over the attempts linked
to each assessment; per-knowledge-component scores are additionally
available through \dbtable{KcEvidence}. From the paired scores the analysis
computes raw gain, Hake's normalized gain
$\langle g \rangle = (\text{post}-\text{pre})/(\text{max}-\text{pre})$
\citep{hake1998}, and an effect size (Cohen's $d$ for paired samples, or a
rank-biserial correlation under the non-parametric route)
\citep{cohen1988}.

Four data limitations shape the method and are handled explicitly. There
is no ``pre''/``post'' marker in the database---the two assessments are
distinguished only by title---so the analysis resolves them by title.
There is no enforced attempt-to-assessment relation, so the join is
by a nullable identifier and must be validated. Both assessments are
practice-mode, permitting multiple attempts, so the analysis must fix a
policy (first, best, or last attempt) and state it. And no maximum-score
or normalized-gain field is stored, so the maximum is summed from item
points and the gain is computed in analysis.

\section{Analyses the data support}
\label{sec:methodology-analyses}

Beyond the two primary questions, the collected signals support a range of
secondary analyses; \cref{tab:analyses} lists ten, ordered roughly by
ascending effort, each with its data basis. They are proposed here and
scaffolded for results in \cref{ch:results}; the report does not report
outcomes for them beyond what the study data yield.

\begin{table}[htbp]
\centering
\small
\caption{Analyses supported by the collected data.}
\label{tab:analyses}
\begin{tabular}{@{}p{0.04\linewidth}p{0.5\linewidth}p{0.38\linewidth}@{}}
\toprule
\# & Analysis & Data basis \\
\midrule
1 & Intervention completion vs.\ dismissal per strategy & \dbtable{LearningIntervention} status by type \\
2 & Confusion$\rightarrow$frustration transition frequency & \dbtable{AffectiveStateWindow} sequences; \dbtable{EfDetection} \\
3 & SM-2 rating calibration vs.\ subsequent recall & \dbtable{SpacedRepetitionCard} $\bowtie$ rating events \\
4 & Pre$\rightarrow$post learning gain per student and KC & \dbtable{Attempt}/\dbtable{GradingResult}/\dbtable{KcEvidence} \\
5 & Cognitive-load index vs.\ assessment performance & \code{derived\_cognitive\_load} $\bowtie$ scores \\
6 & Self-report vs.\ detected affect (validation) & \dbtable{ActivityLog} self-report $\bowtie$ \dbtable{AffectiveStateWindow} \\
7 & Attention-allocation decay over session time & gaze/AOI by session minute \\
8 & Question depth (Bloom's revised taxonomy \citep{anderson2001}) vs.\ performance & \dbtable{Question} tags $\bowtie$ scores \\
9 & At-risk precursors (multimodal) & \code{derived\_at\_risk\_flags} $\bowtie$ preceding-window features \\
10 & Learning velocity vs.\ affective/EF load & \code{derived\_learning\_velocity} $\bowtie$ window aggregates \\
\bottomrule
\end{tabular}
\end{table}

\section{Statistics readiness}
\label{sec:methodology-stats}

The proposal specified validation (item 15), a $t$-test and one-way ANOVA
across phases (item 16), a Mann--Whitney fallback under non-normality
(item 17), and effect sizes (item 18). \Cref{tab:stats-readiness} states,
for each, what the data support and what is missing---the honest core of
the report's statistics story.

\begin{table}[htbp]
\centering
\small
\caption{Statistics readiness for proposal items 15--18.}
\label{tab:stats-readiness}
\begin{tabular}{@{}p{0.06\linewidth}p{0.44\linewidth}p{0.42\linewidth}@{}}
\toprule
Item & Data present & Missing / caveat \\
\midrule
15 & Five-option self-report and windowed machine states with matching
label spaces & Self-report sparse and only post-2026-07-01; no agreement
metric stored---computed in analysis \\
16 & Continuous dependent variables (scores, velocity, cognitive-load
index, detection confidence) & \textbf{No phase/condition marker in the
database}; grouping must be defined externally \\
17 & Same dependent variables, exportable & Same grouping gap; small $n$
threatens power; normality screening is analysis-side \\
18 & All raw quantities for Cohen's $d$, $\eta^2$, rank-biserial & No
stored effect-size or normalized-gain field \\
\bottomrule
\end{tabular}
\end{table}

The single most consequential gap is item 16's: \emph{the database
contains no phase, condition, arm, or group field.} Any phase-1-versus-2
comparison must therefore be defined externally to the data. The report's
position, defended in \cref{ch:results,ch:discussion}, is to operationalize
the two ``phases'' as the pre-study and post-study conditions---the only
principled before/after contrast the data support---and to state that
choice explicitly wherever a grouped test is reported, rather than to
imply a phase design the schema does not encode.

\medskip
\noindent With the methodology fixed, the next chapter presents the study
results within the scaffolding these designs define.
